
On
notera que nous nous écartons ici de la terminologie usuelle (cf. [1]), où on appelait \(\lambda\)-anneau ce que nous appelons pré-\(\lambda\)-anneau, et \(\lambda\)-anneaux « spéciaux » nos \(\lambda\)-anneaux; la raison en est que seuls les \(\lambda\)-anneaux « spéciaux » sont utilisables en pratique.

\textbf{Exemples 2.5.}
\begin{enumerate}[label=\alph*)]
\item Nous avons déterminé toutes les structures de pré-\(\lambda\)-anneau de \(\mathbb Z\)~; il résulte de 2.4.1 que seule celle qui est définie par \(\lambda_t(1)=1+t\) est susceptible d'être une structure de \(\lambda\)-anneau. On a alors \(\lambda_t(n)=(1+t)^n\), d'où~:
\[
  \forall n,n'\in\mathbb Z,\quad
  \lambda_t(n)\circ\lambda_t(n')=(1+t)^n\circ(1+t)^{n'}
  =(1+t)^{nn'}=\lambda_t(nn')~;
\]
\[
  \lambda_u((1+t)^n)
  =
  (\lambda_u(1+t))^n
  =
  (\underline{1}+\{1+t\}u)^n
  =
  \Ghat(\lambda_t)((1+u)^n)
  =
  \Ghat(\lambda_t)(\lambda_u(n)).
\]
Par suite c'est bien une structure de \(\lambda\)-anneau. Si on désigne par \(\binom{X}{k}\) le polynôme à coefficients rationnels \(\frac{1}{k!}X(X-1)\ldots(X-k+1)\), on a donc~:
\begin{equation}\tag{2.5.1}
  \lambda^i(n)=\binom{n}{i}.
\end{equation}

\item Pour tout anneau \(K\), le pré-\(\lambda\)-anneau \(1+K[[t]]^+\) est un \(\lambda\)-anneau. D'abord, \(\lambda_u(\underline{1})=\lambda_u(1+t)=\underline{1}+\{1+t\}u=\underline{1}+\underline{1}.u\). On vérifie les autres relations par la commutativité des diagrammes~:
\[
\begin{tikzcd}[column sep=tiny]
\Ghat^t\times\Ghat^t \arrow[r,"\lambda_u\times\lambda_u"] \arrow[d,"\circ"'] &
(\Ghat^u\circ\Ghat^t)\times\Ghat^u\circ\Ghat^t) \arrow[d,"\circ"]\\
\Ghat^t \arrow[r,"\lambda_u"] &
\Ghat^u\circ\Ghat^t,
\end{tikzcd}
\qquad
\begin{tikzcd}[column sep=tiny]
\Ghat^t \arrow[r,"\lambda_v"] \arrow[d,"\lambda_u"'] &
\Ghat^v\circ\Ghat^t \arrow[d,"\Ghat^v(\lambda_u)"]\\
\Ghat^u\circ\Ghat^t \arrow[r,"\lambda_v"] &
\Ghat^v\circ\Ghat^u\circ\Ghat^t,
\end{tikzcd}
\]
